% Physical interpretation and implementation boundaries added for v1.10.
\section{v1.10 regional stabilization and stationary closure}
\label{sec:v110-physics}

The preceding physical equations and algorithms were copied from the v1.00
document. Numerical examples in those inherited sections describe v1.00
calculations unless explicitly remeasured. The v1.10 additions below describe
the current source. They separate a finite-duration atmospheric stability
condition from a stationary residual condition.

\subsection{Finite-duration regional condition}

Let $q_j(t)$ be a physical-cell quantity, assigned to a lower, transition,
or upper radial region. The monitor partitions accepted physical time into
windows of length $W$. It forms a time average $\overline q_{j,k}$ in window
$k$ by linear interpolation between accepted endpoint states, splitting
intervals exactly at window boundaries. The first complete window fixes the
reference $\overline q_{j,1}$ for that candidate. The check includes the
largest deviation of a window mean from this reference, the change between
neighboring window means, and the span of all window means:
\[
 D_j = \max\!\left\{
 \max_k\frac{|\overline q_{j,k}-\overline q_{j,1}|}{s_j},\;
 \max_{k>1}\frac{|\overline q_{j,k}-\overline q_{j,k-1}|}{s_j},\;
 \frac{\max_k\overline q_{j,k}-\min_k\overline q_{j,k}}{s_j}
 \right\}.
\]
The separate full-range check is
\[
 A_j=\frac{\max_{t\in\mathrm{candidate}}q_j(t)
             -\min_{t\in\mathrm{candidate}}q_j(t)}{s_j}.
\]
These tests are applied to density, velocity, pressure, temperature, electron
and total particle number density, and the atomic species populations. Positive
fluid quantities and regional inventories use their first-window means as
scales; velocity uses an initial sound-speed scale, and species populations
use an initial hydrogen-nucleus density scale. Therefore the numerical limits
are not relative error bounds on a rare population or a near-zero velocity.

The regional inventories are total mass and \emph{internal} energy over the
physical cells of each region; internal energy subtracts kinetic energy from
the conserved total energy. The mass-flow quantity is the actual numerical
face mass flux at the upper face of each region, with the accepted Runge--Kutta
stage weights $1/6$, $1/6$, and $2/3$. Its window-mean change and span are
checked, while its instantaneous range is recorded but not used as a stop
condition. A zero reference flow uses an acoustic-flow scale at machine
precision to keep the denominator nonzero.

The proposed defaults are $D_j\leq0.01$ and $A_j\leq0.05$ for upper cell
quantities, and $D_j\leq0.10$, $A_j\leq0.20$ for lower and transition cell
quantities. Both inventory measures and mass-flow means have limits of 0.01
above the upper cut and 0.05 below it. Every active condition must hold for
eight consecutive windows containing at least 40 accepted endpoints each.
The transition region may be empty, but the lower and upper regions must be
nonempty. An unsuccessful candidate resets its fixed reference. A settling
interval, if specified, is excluded from the first candidate's extrema.

The automatic window is twice the initial outgoing acoustic travel time.
The default maximum duration is the settling interval plus twice the full
confirmation duration. The accepted endpoint at a window edge is included in
the last range test. Rejected numerical steps contribute neither time nor
states. The source reports the effective cuts, limits, sample counts, and
duration in \texttt{regional\_convergence\_status.txt}.

This finite-duration result supports the statement that the modeled upper
atmosphere changed by less than the stated limits while lower-atmosphere
variations remained within their stated limits over the tested duration. It
does not prove stationarity, infinite-time boundedness, or observational
accuracy. A separate stationary residual calculation is required for a
stationary claim.

\subsection{Lower-boundary thermochemical closure in the stationary problem}

The stationary residual in v1.10 evaluates a thermochemical state for the
lower ghost cells at a fixed physical column. The same boundary calculation
is queried when the stored state is certified. The routines
\texttt{stationary\_ghost\_state} and
\texttt{boundary\_chemistry\_of\_state} are defined in
\texttt{stationary\_boundary\_chemistry.f90}; the reservoir composition is
formed in \texttt{reservoir\_composition.f90}. A boundary closure that fails
its own convergence test cannot be treated as a stationary root, even if a
different residual or profile diagnostic is small. The enabled fluid,
composition, radiation, and boundary conditions must all be assessed at the
same state and under the stated input conditions.

The regional test is evaluated in global physical marching of an atomic
atmosphere. It does not evaluate the stationary boundary closure or certify a
Newton--Krylov result, and no automatic transition between the two conditions
is implemented. Existing Huang and Salz results have not been newly accepted
by the definitions or thresholds in this section.
